\chapter{潜在问题、技术难点与改进建议}

\section{风险分级总览}

\begin{longtable}{p{0.09\textwidth}p{0.28\textwidth}p{0.50\textwidth}}
  \toprule
  优先级 & 问题 & 影响与建议\\
  \midrule
  \endfirsthead
  \toprule
  优先级 & 问题 & 影响与建议\\
  \midrule
  \endhead
  P0 & 版本与产物不可严格追溯 & tooth-47 输出早于当前源码/YAML，且文件名显示模式漂移。冻结版本后重跑，清单中写入 Git 提交、全部输入/配置哈希、Blender 和依赖版本。\\
  P0 & 医疗/制造验收边界不完整 & 当前 QA 是几何与拓扑检查，不等同于临床安全、打印可制造性或试戴。必须保留人工审核与实体制造验证。\\
  P1 & 手机避让未被独立复算 & 生成阶段执行差集，验证器不重建包络。增加包络净空探针、关键姿态碰撞检查和梁连续性专项报告。\\
  P1 & 观察窗失败关闭语义不完全一致 & 适配层在 QA 非全真时会记录 \texttt{retained\_after\_failed\_QA}，随后仍可能返回 cutter；与文档中“最终失败关闭”的表述存在差异。生产路径应显式拒绝失败 QA。\\
  P1 & input 模式仍经过体素融合 & “保留输入导管”是相对于不重钻孔而言，并非逐三角面原样保留。增加重采样前后 Hausdorff 距离、孔径和 C 口面积检查。\\
  P1 & 外部工作区模块未在包依赖中表达 & 牙位识别和观察窗依赖 \projectpath{tooth_guide_mapping} 与 \projectpath{scripts}，但 \texttt{pyproject.toml} 只声明 PyYAML。建议打包成显式依赖或建立版本锁定接口。\\
  P2 & 文档与病例参数漂移 & 文档写 0.2 mm/90° 与 0.5/1/2 mm 回退，tooth-47 实际为 0.5 mm 与 1/2/3 mm。由配置自动生成参数表，减少手写常数。\\
  P2 & 部分阶段仍标记 experimental & 只有导管阶段为 stable，其他阶段版本为 0.1--0.4。对每阶段建立明确稳定条件和回归病例覆盖矩阵。\\
  P2 & 操作特征识别假设较强 & 若非导管组件中没有近圆形候选，\texttt{min(circular)} 会失败；圆度阈值和“最近双导管中点”对异常装配敏感。增加清晰诊断与显式配置回退。\\
  P2 & 高密度手机包络成本很高 & tooth-47 包络超过 318 万面。考虑分层包围体、角度自适应采样、SDF/体素并集或只在导板邻域裁剪。\\
  P3 & 输出目录清理范围粗 & 每次生成会删除输出目录顶层全部 PNG/STL。保持专用运行目录，或按清单只清理本程序上次生成的文件。\\
  \bottomrule
\end{longtable}

\section{几何建模本身的难点}

\subsection{薄壳、共面和非流形网格}

牙列和导板来自扫描或上游 CAD，可能存在重复面、孔洞、自交和巨大三角形。薄壳差集尤其容易
在共面处产生碎片。项目通过体素融合提高鲁棒性，再用 manifold3d 精确复切恢复功能孔；
这是一种务实折中，但需要持续量化体素误差。

\subsection{语义方向比几何位置更难}

同一个世界坐标方向不能同时代表所有牙位的“外侧”。牙弓弯曲使局部 outward、近远中和
切向随牙位变化；如果 FDI 左右翻转，观察窗、主梁端点和特殊末端结构会一起翻转。相比
布尔失败，这类错误更危险，因为结果可能仍然是漂亮的封闭 STL。

\subsection{局部最优与整体结构}

每个锚点都可以局部可行，但四根主梁和 Y 梁组合后仍可能过近、交叉或干涉手机包络。当前
流程主要以确定性规则和少量 maximin/最短距离选择构造结果，尚未形成统一的多目标优化：
例如同时最小化梁长、曲率和体积，最大化夹角、根部面积和器械净空。

\subsection{布尔顺序是功能语义的一部分}

generated 和 input 模式的差异不是实现细节。一个错误的“最后再统一复切”会改变真实输入
导管；一个错误的“先切一次就结束”会让体素融合堵回孔道。后续重构必须用模式化测试锁定
每一步 cutter 的作用对象和顺序。

\section{建议的改进路线}

\begin{enumerate}
  \item \term{先补追溯}：每次生成写出统一 run manifest，包括源码版本、dirty 状态、输入和
  配置哈希、算法版本、命令、耗时、环境和输出哈希。
  \item \term{再补验证空白}：优先增加手机包络独立净空、操作窗终态检查，以及 input 导管
  重采样误差和真实孔径量测。
  \item \term{收紧失败语义}：任何观察窗 QA 非全真都不允许进入正式生成；警告继续的贴合脚
  应在生产模式转为明确的人工签核项。
  \item \term{建立病例覆盖矩阵}：普通、断裂导板、U 型末端、远中公共节点、相邻双种植位、
  input/generated 两模式均有固定回归与阈值。
  \item \term{控制性能}：为手机包络和观察窗记录每阶段耗时、峰值内存和网格规模，避免只在
  最终命令超时后排查。
  \item \term{自动生成文档参数}：从 dataclass 默认值和病例 YAML/JSON 生成表格，降低说明文档
  与实际配置分叉。
\end{enumerate}

